\section{Recomendaciones}
\label{sec:recomendaciones}

[Formular las recomendaciones tecnicas y operativas derivadas del analisis.
Cada recomendacion debe estar directamente sustentada en una conclusion o hallazgo
de la seccion anterior. Clasificar por categoria (ingenieria, operacional,
infraestructura, calificacion, mantenimiento, etc.).]

\begin{table}[htbp]
    \centering
    \caption{[Leyenda descriptiva de las recomendaciones.]}
    \label{tab:recomendaciones}
    \setcounter{itemcount}{0}
    \begin{tabularx}{\textwidth}{@{}NlX@{}}
        \toprule
        \multicolumn{1}{c}{\textbf{Item}} & \textbf{Categoria} & \textbf{Recomendacion} \\
        \midrule
        & [Categoria] & [Descripcion de la recomendacion tecnica.] \\
        & [Categoria] & [Descripcion de la recomendacion tecnica.] \\
        & [Categoria] & [Descripcion de la recomendacion tecnica.] \\
        \bottomrule
    \end{tabularx}
\end{table}

\noindent\textbf{[Titulo de recomendaciones adicionales.]}

\begin{table}[htbp]
    \centering
    \caption{[Leyenda descriptiva de recomendaciones complementarias.]}
    \label{tab:recomendaciones_add}
    \setcounter{itemcount}{0}
    \begin{tabularx}{\textwidth}{@{}NlX@{}}
        \toprule
        \multicolumn{1}{c}{\textbf{Item}} & \textbf{Categoria} & \textbf{Recomendacion} \\
        \midrule
        & [Categoria] & [Descripcion de la recomendacion complementaria.] \\
        & [Categoria] & [Descripcion de la recomendacion complementaria.] \\
        \bottomrule
    \end{tabularx}
\end{table}
